\begin{textsample}{POS Dim 2 – 2011\_c4 – Score 9.00 – 2011\_c4/enade\_2011\_eng\_g\_6\_7d2177b9ec.txt}  \label{ex:f2_pos_130}
CAPÍTULO 4 PERCEPÇÃO DA PROVA As análises feitas neste capítulo tratam das percepções dos concluintes da Área de Engenharia – Grupo VI sobre a prova aplicada no ENADE/2011.
Estas percepções foram mensuradas por meio de nove questões que avaliaram desde o grau de dificuldade da prova até o tempo gasto para concluí -la.
As percepções sobre a prova foram relacionadas com o desempenho dos estudantes e com a Grande Região de funcionamento do curso.
O \textbf{questionário} de percepção da prova encontra -se ao final do Anexo IV, que traz a \textbf{reprodução} da prova.
O desempenho dos estudantes foi classificado em quatro quartos.
Para tanto, esse desempenho foi \textbf{ordenado} de forma ascendente.
O percentil 25, P25, também conhecido como primeiro quartil, é a nota de desempenho que deixa um quarto ( 25 \% ) dos valores observados abaixo e três quartos acima.
A Figura 1 apresenta uma ilustração deste conceito.
O quarto inferior de desempenho é composto pelas notas abaixo do primeiro quartil.
Já o percentil 75, P75, também conhecido como terceiro quartil, é o valor para o qual há três quartos ( 75 \% ) dos dados abaixo e um quarto acima dele.
O quarto superior de desempenho é composto pelas notas iguais ou acima do terceiro quartil.
O percentil 50, P50, também conhecido como mediana, é o valor que divide as notas em dois conjuntos de igual tamanho.
O segundo quarto inclui valores entre o primeiro quartil ( P25 ) e a mediana.
O terceiro quarto contém os valores entre a mediana ( P50 ) e o terceiro quartil ( P75 ).
Vale ressaltar que percentis, quartis e medianas são pontos que não obrigatoriamente pertencem ao conjunto original de dados, ao passo que os quartos são \textbf{subconjuntos} dos dados originais.
A seguir, serão apresentados gráficos com resultados selecionados, relativos às nove questões avaliadas por grupos de estudantes.
Os gráficos apresentam nas barras o percentual de alunos que assinalaram uma das opções ou a soma das porcentagens daqueles que assinalaram duas ( ou três ) delas.
Por exemplo, para as questões 1 e 2, os gráficos apresentam a porcentagem total de participantes que assinalaram as opções ( D ) difícil e ( E ) muito difícil.
Em cada barra foram assinalados também os extremos do intervalo de confiança de 95 \% como linhas verticais unidas por uma linha horizontal na forma da letra H maiúscula, semelhantemente aos gráficos do Capítulo 3.
As Tabelas no Anexo II apresentam os valores absolutos e a distribuição percentual das alternativas válidas das nove questões, segundo o mesmo \textbf{recorte} de desempenho dos alunos e Grande Região de funcionamento do curso. 4.1 GRAU DE DIFICULDADE DA PROVA 4.1.1 Componente de Formação Geral Ao avaliarem “ Qual o grau de dificuldade desta prova na parte de Formação Geral? ” ( Questão 1 ), 21,8 \% do grupo de inscritos e presentes optaram pelas alternativas difícil ou muito difícil.
Entretanto, para mais da metade dos estudantes ( 60,5 \% ), o Componente de Formação Geral da prova foi considerado com grau de dificuldade médio ( Gráfico 4.1, Gráfico 4.2 e, no Anexo II, a Tabela II.1 ).
O percentual de estudantes que consideraram a prova como difícil ou muito difícil foi maior na região Sul, onde a proporção foi de 26,0 \%, enquanto a de menor incidência foi a Nordeste, com 15,6 \%.
No Gráfico 4.1 é possível observar que esta diferença é estatisticamente significativa.
Nas Grandes Regiões, a proporção de presentes à prova que consideraram o Componente de Formação Geral como sendo de grau de dificuldade médio esteve entre 57,3 \% na região Sul e 65,7 \% na região Centro-Oeste.
O percentual de alunos que consideraram a prova difícil ou muito difícil apresentou tendência decrescente conforme o aumento do desempenho, variando de 28,7 \% no 1º quarto até 13,7 \% no 4º quarto.
A alternativa modal para esta pergunta foi médio, com tendência crescente entre o 1º ( 59,7 \% ) e o 3º ( 63,7 \% ) quarto, declinando o valor percentual no 4º quarto ( 58,5 \% ). 4.1.2 Componente de Conhecimento Específico Ao responderem à Questão 2 – “ Qual o grau de dificuldade desta prova na parte de Componente Específico? ” – 43,5 \% do grupo de estudantes classificaram -na como difícil ou muito difícil.
Além disso, o Componente de Conhecimento Específico da prova foi considerado com grau de dificuldade médio por 51,7 \% dos alunos ( Gráfico 4.3, Gráfico 4.4, e, no Anexo II, a Tabela II.2 ).
A análise das respostas dos estudantes quanto ao grau de dificuldade do Componente de Conhecimento Específico da prova, agregado por Grande Região, mostra que a diferença entre a maior e a menor proporção de alunos que a avaliaram como difícil ou muito difícil é estatisticamente significativa : 36,2 \% na região Nordeste e 51,5 \% na Sul.
O percentual de alunos que classificaram o grau de dificuldade como médio, no Componente de Conhecimento Específico, variou de 44,6 \% a 59,2 \%, para as mesmas regiões, Sul e Nordeste, respectivamente.
Considerando -se a avaliação da dificuldade das questões do Componente de Conhecimento Específico da prova, de acordo com o desempenho dos estudantes, não se observa diferença estatisticamente significativa.
Em todos os quartos, a proporção dos que classificaram a parte específica como difícil ou muito difícil variou de 42,4 \% ( 1º quarto ) a 45,4 \% ( 2º quarto ).
A alternativa modal para a Questão 2 foi o médio, com 50,9 \% do quarto inferior e 52,1 \% do superior optando por esta resposta. 4.2 EXTENSÃO DA PROVA EM RELAÇÃO AO TEMPO TOTAL \textbf{Indagados} quanto à extensão da prova, em relação ao tempo total oferecido para a sua resolução ( Questão 3 ), os estudantes apontaram, com maior incidência, a alternativa que considerava a extensão adequada ( Gráfico 4.5, Gráfico 4.6, e, no Anexo II, a Tabela II.3 ).
O percentual de alunos que responderam ser a extensão da prova adequada foi de 42,6 \%.
Já 51,8 \% dos inscritos presentes consideraram que o exame foi longo ou muito longo e menos do que 6 \% o avaliaram como curto ou muito curto.
Entre as Grandes Regiões a proporção daqueles que avaliaram a prova como longa ou muito longa em relação ao tempo total destinado à sua resolução variou pouco : de 50,4 \% na região Sul até 60,9 \% na região Nordeste.
A diferença entre as regiões Sul e Nordeste é estatisticamente significativa.
Considerando -se o desempenho dos alunos, a proporção dos que optaram pela extensão da prova adequada apresentou uma tendência decrescente com o desempenho do 1º quarto ( 45,0 \% ) ao 3º quarto ( 40,3 \% ), ocorrendo um aumento do valor percentual do 3º para o último quarto ( 41,6 \% ).
No Gráfico 4.6, pode -se constatar que para os quatro quartos de desempenho a proporção de estudantes que consideraram a prova longa ou muito longa, em relação ao tempo total destinado à sua resolução, manteve -se em torno de 52 \%, sem diferenças estatisticamente significativas. 4.3 COMPREENSÃO DOS ENUNCIADOS DAS QUESTÕES 4.3.1 Componente de Formação Geral Com relação aos enunciados das questões do Componente de Formação Geral ( Questão 4 ), as opiniões foram positivas, já que 73,9 \% dos alunos avaliados consideraram os enunciados de todas ou da maioria das questões claros e objetivos ( Gráfico 4.7, Gráfico 4.8, e, no Anexo II, a Tabela II.4 ).
Na análise regional, a \textbf{percentagem} de estudantes que avaliaram que todos ou a maioria dos enunciados das questões do Componente de Formação Geral estavam claros e objetivos variou de 70,8 \% na região Sul a 78,2 \% na região Norte.
A única diferença estatisticamente significativa entre as Grandes Regiões ocorreu entre as regiões Sul e Nordeste ( 77,0 \% ) A análise das percepções dos estudantes sobre a clareza e objetividade dos enunciados permite afirmar que todos ou a maioria dos enunciados de questões relativas ao Componente de Formação Geral foram considerados claros e objetivos para a maior parte dos respondentes ( maior do que 70 \% em todas as regiões e maior do que 65 \% para todos os quartos de desempenho ).
Segundo o desempenho, observa -se que a proporção dos que emitiram esta opinião \textbf{cresce} conforme o desempenho aumenta, com diferenças estatisticamente significativas.
No quarto superior, a clareza e objetividade de todos ou da maioria dos enunciados das questões foi percebida por 81,6 \% dos alunos e no quarto de desempenho inferior tal avaliação foi emitida por 65,1 \% deles, numa proporção crescente com o desempenho. 4.3.2 Componente de Conhecimento Específico Com relação aos enunciados das questões do Componente de Conhecimento Específico da prova, para 69,4 \% dos estudantes avaliados da Área de Engenharia – Grupo VI a clareza e a objetividade ( Questão 5 ) estavam presentes em todas ou na maioria das questões ( Gráfico 4.9, Gráfico 4.10, e no Anexo II, a Tabela II.5 ).
A maioria dos estudantes de todas as Grandes Regiões brasileiras considerou claros e objetivos todos ou a maioria dos enunciados das questões do Componente de Conhecimento Específico da prova, percentual sempre maior do que 65 \%.
A diferença entre 66,7 \% ( Centro-Oeste ) e 74,8 \% ( Norte ) não é estatisticamente significativa.
A proporção de estudantes que consideraram os enunciados das questões claros e objetivos apresenta uma tendência crescente em relação ao aumento de desempenho : mais elevada no quarto superior ( 73,4 \% ) se comparada ao quarto inferior de desempenho ( 63,7 \% ).
Esta diferença é estatisticamente significativa. 4.4 SUFICIÊNCIA DAS INFORMAÇÕES/INSTRUÇÕES FORNECIDAS Ao avaliarem as informações/instruções fornecidas para a resolução das questões ( Questão 6 ), 85,9 \% dos respondentes da Área de Engenharia – Grupo VI de todo o Brasil afirmaram que estas eram até excessivas ou suficientes em todas ou na maioria das questões ( Gráfico 4.11, Gráfico 4.12, e, no Anexo II, a Tabela II.6 ).
Quanto à distribuição de respondentes pelas Grandes Regiões observa -se que a proporção de estudantes que consideraram as informações/instruções fornecidas até excessivas ou suficientes em todas ou na maioria das questões foi sempre superior a 75 \%, chegando a 86,4 \% na região Sul.
As diferenças entre as regiões não chegaram a ser estatisticamente significativas.
Levando -se em conta o desempenho dos participantes, notam -se diferenças estatisticamente significativas entre as opiniões de estudantes do quarto inferior e superior de desempenho, como mostra o Gráfico 4.12.
O percentual de participantes que avaliaram as informações/instruções como até excessivas ou suficientes em todas ou na maioria das questões foi mais elevado no quarto superior ( 91,3 \% ), percentual superior à média nacional ( 85,9 \% ).
Já no quarto inferior, a suficiência das informações/instruções declarada como até excessiva em todas ou na maioria das questões foi percebida por 81,2 \% dos respondentes.
Nota -se uma proporção crescente com o desempenho na escolha destas alternativas em conjunto. 4.5 DIFICULDADE ENCONTRADA AO RESPONDER À PROVA Perguntados sobre as dificuldades com as quais se depararam ao responder à prova ( Questão 7 ), 9,2 \% dos estudantes apontaram o desconhecimento do conteúdo.
Para 45,5 \%, a forma diferente de abordagem do conteúdo foi indicada como dificuldade.
Já a falta de motivação para fazer a prova foi a dificuldade apontada por 25,6 \% dos respondentes.
Considerando -se todo o Brasil, 14,4 \% dos respondentes afirmaram que não tiveram qualquer tipo de dificuldade para responder à prova ( Tabela II.7 no Anexo II ).
Os Gráficos 4.13 e 4.14 apresentam os percentuais de estudantes que apontaram o desconhecimento do conteúdo como dificuldade percebida ao responder à prova.
Na análise por Grandes Regiões, o percentual de inscritos e presentes que apontaram o desconhecimento do conteúdo como dificuldade ao responder à prova não superou 14 \%.
Os percentuais variaram de 8,0 \% na região Sudeste a 13,3 \% na Centro-Oeste, não sendo esta diferença significativa estatisticamente.
A forma diferente de abordagem do conteúdo foi a escolha modal dos estudantes, com percentuais que variaram de 44,8 \% ( região Sul ) a 48,2 \% ( Norte ).
O percentual de alunos que citou a falta de motivação como dificuldade variou de 17,3 \% ( região Centro-Oeste ) a 26,7 \% ( região Sudeste ).
Os que declararam não ter qualquer dificuldade para responder à prova variaram de 11,3 \% na região Nordeste a 15,2 \% na Sudeste.
Com relação aos quartos de desempenho, os valores percentuais extremos, para a opção desconhecimento do conteúdo ocorreu no 2º quarto ( 8,8 \% ) e no 4º quarto ( 10,1 \% ).
A alternativa modal para os alunos, quando agregados pelos quartos de desempenho, foi que a dificuldade encontrada foi causada pela forma diferente de abordagem do conteúdo, ocorrendo os valores percentuais extremos no 4º quarto ( 41,3 \% ) e no 2º quarto ( 47,8 \% ). 4.6 CONTEÚDOS DAS QUESTÕES OBJETIVAS DA PROVA Ao analisarem os conteúdos das questões objetivas da prova ( Questão 8 ), um percentual muito pequeno dos estudantes avaliados, apenas 3,2 \%, afirmou que não estudou ainda a maioria desses conteúdos ( Gráficos 4.15, Gráfico 4.16, e a Tabelas II.8 no Anexo II ).
A maioria ( 73,7 \% ) afirmou ter estudado e aprendido muitos ou todos os conteúdos.
Na análise por Grande Região, a proporção de respondentes que escolheram a opção não estudou ainda a maioria desses conteúdos, foi pequena.
Observa -se que nas regiões Norte ( 3,5 \% ), Nordeste ( 3,8 \% ), Sul ( 4,5 \% ) e Centro-Oeste ( 4,1 \% ), apesar de pequenas, as proporções foram maiores do que a média nacional ( 3,2 \% ).
Não se observa diferença estatisticamente significativa entre as regiões.
Em todas as regiões, a maioria dos presentes afirmou ter estudado e aprendido muitos ou todos os conteúdos, com proporções variando entre 66,5 \% na região Norte e 75,5 \% na Sudeste.
Considerando -se separadamente as opiniões de estudantes dos quatro quartos de desempenho, observa -se que, no quarto inferior, 5,2 \% ofereceram como resposta que não estudou ainda a maioria desses conteúdos, sendo menos de 1,5 \% os do quarto superior com a mesma resposta.
A diferença entre os alunos que optaram por este motivo de dificuldade nos quartos extremos é estatisticamente significativa e a proporção se mostra decrescente com o desempenho.
Tendo em conta o quarto superior, 81,5 \% dos alunos afirmaram ter estudado e aprendido muitos ou todos os conteúdos. 4.7 TEMPO GASTO PARA CONCLUIR A PROVA Ao responderem sobre o tempo de conclusão da prova ( Questão 9 ), mais de dois terços dos estudantes ( 71,8 \% ) afirmaram ter gasto entre duas e quatro horas ( Gráfico 4.17, Gráfico 4.18 e, no Anexo II, a Tabela II.9 ).
Considerando -se as cinco Grandes Regiões brasileiras, os que utilizaram entre duas e quatro horas para finalizar a prova superaram o percentual nacional na região Sudeste ( 72,2 \% ).
Em três das cinco Grandes Regiões, Norte, Nordeste e Sul, o percentual de alunos que \textbf{dispensaram} entre duas e quatro horas para concluir a prova, ficou em torno de 71 \% e na região Centro-Oeste o percentual ficou muito abaixo das demais com 65,7 \%, como mostra o Gráfico 4.17.
Uma vez sendo analisadas as alternativas escolhidas pelos estudantes que se situam nos diferentes quartos de desempenho, observa -se que uma maior proporção de participantes no quarto superior declarou ter gasto entre duas e quatro horas para concluir a prova quando comparados com os do quarto inferior, respectivamente 78,2 \% e 65,2 \%.
As diferenças entre alguns dos quartos são estatisticamente significativas e os valores evidenciam uma tendência crescente com o desempenho.

% matched lemmas: crescer, dispensar, indagar, ordenar, percentagem, questionário, recorte, reprodução, subconjunto
\end{textsample}
